\documentclass[conference,a4paper]{IEEEtran}

\usepackage[utf8]{inputenc}
\usepackage[T5]{fontenc}
\usepackage[vietnamese,english]{babel}
\usepackage{amsmath}
\usepackage{amssymb}
\usepackage{booktabs}
\usepackage{array}
\usepackage[protrusion=false,expansion=true]{microtype}
\usepackage[unicode]{hyperref}
\usepackage{bookmark}
\usepackage[table]{xcolor}
\usepackage{tikz}
\usepackage{graphicx}
\usepackage{balance}
\usetikzlibrary{arrows.meta,positioning,shapes.geometric}

\hypersetup{colorlinks=true,linkcolor=black,citecolor=blue,urlcolor=blue,pdftitle={Driver Drowsiness and Reduced Attention Detection Using a Multimodal In-Cabin Model}}
\addto\captionsvietnamese{\renewcommand{\refname}{References}}
\renewcommand{\arraystretch}{1.08}
\setlength{\emergencystretch}{1.5em}
\newcolumntype{L}[1]{>{\raggedright\arraybackslash}p{#1}}

\title{\fontsize{20}{22}\selectfont PHÁT HIỆN BUỒN NGỦ VÀ GIẢM CHÚ Ý CỦA TÀI XẾ BẰNG MÔ HÌNH ĐA PHƯƠNG THỨC TRONG CABIN}

\author{
\IEEEauthorblockN{Ngo Van Minh, Nguyen Thi Phuong}
\IEEEauthorblockA{Group 01\\
Dai Nam University\\
Hanoi, Vietnam}
}

\begin{document}
\maketitle

\begin{abstract}
This paper studies multimodal driver-state recognition in an in-cabin setting. The core model, \textit{mmrt\_multimodal\_full\_v2\_t4x2}, is trained on the \textit{DBR-DMD-v1.1 CLEAN\_FULLMM\_5CLASS} dataset with four synchronized modalities: body, face, audio, and skeleton. The best verified checkpoint reaches a validation macro-F1 of 0.8066. Artifact-based analysis shows stable recognition for \textit{phone\_use}, \textit{eye\_closure}, and \textit{yawning}, while the dominant confusion remains between \textit{normal\_driving} and \textit{looking\_away}. These results support temporal multimodal fusion as a practical direction for driver safety assessment when a single visual cue is insufficient. The \textit{MDDAD} system is discussed only as an application layer for exposing the verified model through upload analysis, monitoring, and reporting workflows.
\end{abstract}

\begin{IEEEkeywords}
driver safety, multimodal learning, in-cabin monitoring, drowsiness detection, distraction detection, risk scoring
\end{IEEEkeywords}

\selectlanguage{vietnamese}

\noindent\textbf{\textit{Tóm tắt}}---Bài báo nghiên cứu bài toán nhận diện trạng thái buồn ngủ và suy giảm chú ý của tài xế trong cabin theo hướng đa phương thức. Mô hình trung tâm là \textit{mmrt\_multimodal\_full\_v2\_t4x2}, được huấn luyện trên bộ dữ liệu \textit{DBR-DMD-v1.1 CLEAN\_FULLMM\_5CLASS} với bốn nguồn tín hiệu đồng bộ gồm \textit{body}, \textit{face}, \textit{audio} và \textit{skeleton}. Checkpoint tốt nhất đạt \textit{macro-F1 = 0.8066} trên tập validation. Phân tích artifact cho thấy mô hình nhận diện ổn định các lớp \textit{phone\_use}, \textit{eye\_closure} và \textit{yawning}, trong khi sai số chính tập trung ở ranh giới giữa \textit{normal\_driving} và \textit{looking\_away}. Kết quả này cho thấy hợp nhất tín hiệu đa phương thức theo thời gian là một hướng khả thi cho giám sát an toàn tài xế khi tín hiệu thị giác đơn lẻ không đủ mô tả hành vi. Hệ thống \textit{MDDAD} chỉ được sử dụng như lớp ứng dụng hóa để biểu diễn mô hình qua các workflow phân tích video, giám sát và báo cáo.

\vspace{0.75em}

\section[Gioi Thieu]{Giới Thiệu}
Giám sát an toàn tài xế trong cabin là một bài toán có ý nghĩa trực tiếp với an toàn giao thông, do nhiều hành vi nguy hiểm như buồn ngủ, mất tập trung, nhìn lệch hướng kéo dài, hoặc sử dụng điện thoại thường xuất hiện trong các khoảng thời gian ngắn nhưng có hậu quả lớn. Trong nhiều hệ thống sớm, quyết định thường dựa trên một kênh đầu vào duy nhất hoặc trên các ngưỡng heuristic. Cách tiếp cận này có thể hoạt động ở những tình huống rõ ràng, nhưng thường khó mở rộng khi hành vi được biểu hiện đồng thời qua khuôn mặt, tư thế cơ thể, động học skeleton và tín hiệu âm thanh.

Trong bối cảnh đó, các hướng video backbone mạnh như SlowFast~\cite{slowfast}, TimeSformer~\cite{timesformer} hay Video Swin Transformer~\cite{videoswin} cho thấy tiềm năng biểu diễn cao cho chuỗi hình ảnh, nhưng bài toán trong cabin không chỉ nằm ở video branch. Thách thức trung tâm là cách đồng bộ và hợp nhất nhiều nguồn tín hiệu dị biệt nhau trong cùng một hệ suy luận theo thời gian. Từ góc nhìn này, đề tài không coi multimodal learning là phần mở rộng phụ của một model ảnh tĩnh, mà là trục chính để biến bài toán thành nhận diện hành vi tài xế theo clip, có ngữ cảnh và có khả năng chuyển giao sang hệ thống demo.

Ý tưởng chính của báo cáo là tách rõ hai tầng. Tầng nghiên cứu tập trung vào baseline đa phương thức thống nhất \textit{mmrt\_multimodal\_full\_v2\_t4x2} trên bộ dữ liệu 5 lớp. Tầng thứ hai là \textit{MDDAD}, một application layer biến đầu ra mô hình thành workflow phân tích video tải lên, giám sát webcam, chấm điểm rủi ro và sinh báo cáo. Việc phân tách này quan trọng vì nó giúp báo cáo trung thực với trạng thái triển khai thực tế: baseline nghiên cứu đã được kiểm chứng trên dataset chuẩn, trong khi runtime của app hiện vẫn mang tính lai ghép để ưu tiên độ ổn định demo.

Trong phạm vi báo cáo, các đóng góp chính được chốt lại theo hướng học thuật thay vì theo danh sách tính năng:
\begin{enumerate}
\item Xác lập một baseline nhận diện hành vi tài xế đa phương thức thống nhất trên bộ dữ liệu \textit{DBR-DMD-v1.1 CLEAN\_FULLMM\_5CLASS} với bốn modality \textit{body}, \textit{face}, \textit{audio} và \textit{skeleton}, thay cho cách nhìn đơn kênh hoặc heuristic thuần túy.
\item Chuẩn hóa cách đọc bằng chứng thực nghiệm của bundle \textit{mmrt\_multimodal\_full\_v2\_t4x2}, bao gồm metric checkpoint, input contract, confusion pattern, và giới hạn tái lập; nhờ đó báo cáo chỉ khẳng định những gì đã xác minh trực tiếp.
\item Chỉ ra điểm nghẽn quyết định của baseline hiện tại nằm ở biên \textit{normal\_driving}/\textit{looking\_away}, thay vì ở các lớp có tín hiệu hình thái rõ như \textit{phone\_use}, \textit{eye\_closure} và \textit{yawning}.
\item Tách rời tầng nghiên cứu và tầng ứng dụng hóa \textit{MDDAD} để tránh lẫn giữa đóng góp mô hình và đóng góp triển khai demo.
\end{enumerate}

\section[Nghien Cuu Lien Quan]{Nghiên Cứu Liên Quan}
Có thể chia các hướng nghiên cứu liên quan tới bài toán giám sát an toàn tài xế trong cabin thành ba nhóm. Nhóm thứ nhất là các hệ thống \textit{face-centric} hoặc heuristic-centric, trong đó quyết định chủ yếu dựa trên khuôn mặt, chuyển động mắt, độ mở miệng hoặc góc đầu~\cite{dms,mediapipe}. Đây là hướng rất thực dụng vì chi phí tính toán thấp, dễ tích hợp vào webcam runtime, và cho tín hiệu khá rõ với các hành vi như mắt nhắm kéo dài hoặc ngáp. Tuy nhiên, điểm yếu cơ bản của hướng này là nó coi phần lớn ngữ cảnh còn lại của cabin như nhiễu. Khi người lái chỉ lệch hướng nhìn ngắn, che khuất một phần mặt, hoặc biểu hiện hành vi qua tư thế thân trên nhiều hơn qua gương mặt, quyết định dựa vào một nhánh duy nhất thường trở nên kém ổn định.

Nhóm thứ hai là các mô hình nhận diện hành động dựa trên video, vốn xem hành vi là một vấn đề ngữ cảnh không-thời gian thay vì một tập hợp trạng thái tĩnh. Các backbone như SlowFast~\cite{slowfast}, TimeSformer~\cite{timesformer} và Video Swin~\cite{videoswin} cho thấy khả năng mô hình hóa clip mạnh hơn đáng kể so với các pipeline khung-hình rời rạc. Dẫu vậy, trong bối cảnh cabin, việc chỉ tăng sức mạnh của video backbone vẫn chưa giải quyết trọn vẹn bài toán. Nhiều hành vi nguy cơ không chỉ được biểu lộ qua diện mạo hình ảnh, mà còn qua tín hiệu âm thanh, cấu trúc động học skeleton, hoặc sự phối hợp giữa hướng đầu và thao tác vùng tay. Nói cách khác, chuyển từ ảnh sang video là một bước cần, nhưng chưa đủ nếu bài toán được đặt như một hệ thống đa tín hiệu.

Nhóm thứ ba là các nghiên cứu hành vi tài xế đa phương thức ở mức chi tiết hơn, tiêu biểu như Drive\&Act~\cite{driveact}. Điểm quan trọng rút ra từ hướng này là bài toán trong cabin cần được nhìn như một bài toán \textit{clip-level multimodal reasoning}: đầu vào không chỉ là chuỗi frame, mà là nhiều nguồn tín hiệu đồng bộ phải được hợp nhất trong cùng một quyết định. Đây cũng là điểm mà đề tài hiện tại lựa chọn làm trục chính. Thay vì tiếp tục mở rộng một pipeline heuristic hoặc đẩy toàn bộ hy vọng vào một video backbone lớn, mô hình nghiên cứu được xây dựng như một baseline đa phương thức thống nhất với bốn nhánh \textit{body}, \textit{face}, \textit{audio} và \textit{skeleton}.

\begin{table}[t]
\caption{So sánh các hướng nghiên cứu liên quan và vị trí của đề tài}
\label{tab:related-works}
\centering
\scriptsize
\begin{tabular}{L{1.15cm}L{1.45cm}L{1.0cm}L{1.35cm}L{1.45cm}}
\toprule
Hướng & Công trình đại diện & Đầu vào & Điểm mạnh & Hạn chế chính \\
\midrule
Heuristic khuôn mặt & Fatigue monitoring~\cite{dms}, MediaPipe~\cite{mediapipe} & Landmark & Rẻ, nhanh, hợp demo & Yếu khi che khuất, thiếu ngữ cảnh clip \\
Video recognition đơn kênh & SlowFast~\cite{slowfast}, TimeSformer~\cite{timesformer}, Video Swin~\cite{videoswin} & RGB clip & Mạnh về ngữ cảnh thời gian & Chưa khai thác tốt audio, skeleton, quan hệ đa nguồn \\
Hành vi cabin đa phương thức & Drive\&Act~\cite{driveact} & RGB, IR, depth, pose & Bám sát bản chất hành vi cabin & Khó tái lập đầy đủ, khó chuẩn hóa đầu vào \\
\bottomrule
\end{tabular}
\end{table}

Từ Bảng~\ref{tab:related-works}, có thể thấy khoảng trống mà đề tài nhắm tới không phải là thay thế hoàn toàn các hướng trước đó, mà là nối chúng lại theo một cách phù hợp hơn với bài toán hiện tại. Cụ thể, đề tài kế thừa tính thực dụng của các pipeline nhẹ, kế thừa tư duy suy luận theo clip của các backbone video, và tiến gần hơn tới bản chất hành vi cabin đa nguồn. Vì vậy, mô hình \textit{mmrt\_multimodal\_full\_v2\_t4x2} được định vị không như một kiến trúc ``lớn nhất'' hay ``mới nhất'', mà như một baseline nghiên cứu có phạm vi hợp lý cho bộ dữ liệu hiện có. Cách chọn các branch gọn cho từng modality, sau đó hợp nhất qua \textit{gated fusion} và suy luận chuỗi bằng \textit{BiGRU + attention}, giúp đề tài đạt ba mục tiêu cùng lúc: (1) bám sát bản chất đa nguồn của hành vi tài xế, (2) giữ chi phí mô hình hóa ở mức chấp nhận được với tập dữ liệu 2216 mẫu, và (3) tạo ra một lõi nghiên cứu đủ rõ để có thể phân tích lỗi, so sánh phiên bản và chuyển giao dần sang tầng ứng dụng.

\section[Bai Toan Du Lieu Va Tien Xu Ly]{Bài Toán, Dữ Liệu Và Tiền Xử Lý}
\subsection{Phát biểu bài toán}
Ý tưởng cốt lõi của đề tài là xem giám sát tài xế như một bài toán \textit{phân loại clip đa phương thức theo thời gian}, thay vì nhận diện trạng thái trên từng khung hình rời rạc. Mỗi mẫu trong bộ dữ liệu nghiên cứu là một đơn vị clip có bốn modality đồng bộ: \textit{body}, \textit{face}, \textit{audio} và \textit{skeleton}. Đầu ra của mô hình là một trong năm lớp hành vi: \textit{normal\_driving}, \textit{eye\_closure}, \textit{yawning}, \textit{looking\_away} và \textit{phone\_use}. Cách phát biểu này có hai hệ quả trực tiếp. Thứ nhất, quyết định cuối cùng phải dựa trên bằng chứng tích lũy theo thời gian, không thể chỉ nhìn một frame đơn lẻ. Thứ hai, các modality không đóng vai trò đối xứng; từng nhánh mang các loại tín hiệu khác nhau và chỉ trở nên hữu ích khi được đặt trong cùng một ngữ cảnh clip.

\subsection{Thống kê dữ liệu}
Bộ dữ liệu \textit{DBR-DMD-v1.1 CLEAN\_FULLMM\_5CLASS} gồm 2216 mẫu, chia thành 1432 mẫu huấn luyện, 370 mẫu validation và 414 mẫu kiểm thử. Một điểm quan trọng của bộ dữ liệu này là các mẫu dùng để đánh giá đều có đủ bốn modality, nhờ đó việc so sánh các phiên bản mô hình không bị nhiễu bởi bài toán thiếu kênh đầu vào. Nói cách khác, chênh lệch kết quả giữa các phiên bản chủ yếu phản ánh chất lượng mô hình và chiến lược huấn luyện, thay vì phản ánh sự không đồng nhất của dữ liệu đầu vào.

\begin{table}[t]
\caption{Thống kê dữ liệu theo tập và theo lớp}
\label{tab:data-stats}
\centering
\small
\begin{tabular}{lrrrrr}
\toprule
Tập & normal & eye & yawn & away & phone \\
\midrule
Train & 720 & 80 & 30 & 298 & 304 \\
Val & 181 & 23 & 8 & 83 & 75 \\
Test & 168 & 36 & 16 & 91 & 103 \\
\bottomrule
\end{tabular}
\end{table}

Thống kê ở Bảng~\ref{tab:data-stats} cho thấy bài toán không chỉ có quy mô vừa phải mà còn có mất cân bằng lớp. Hệ quả là accuracy được dùng để tham khảo, nhưng không đủ để làm trục chính khi chọn checkpoint hoặc phân tích chất lượng mô hình.

\begin{table}[t]
\caption{Phân bố lớp toàn bộ bộ dữ liệu và mức mất cân bằng}
\label{tab:class-imbalance}
\centering
\small
\begin{tabular}{lrr}
\toprule
Lớp & Số mẫu & Tỷ lệ (\%) \\
\midrule
\textit{normal\_driving} & 1069 & 48.2 \\
\textit{eye\_closure} & 139 & 6.3 \\
\textit{yawning} & 54 & 2.4 \\
\textit{looking\_away} & 472 & 21.3 \\
\textit{phone\_use} & 482 & 21.8 \\
\bottomrule
\end{tabular}
\end{table}

Bảng~\ref{tab:class-imbalance} làm rõ vì sao \textit{macro-F1} phải được ưu tiên hơn accuracy. Lớp \textit{normal\_driving} chiếm gần một nửa bộ dữ liệu, trong khi \textit{yawning} chỉ chiếm 2.4\%. Nếu chỉ tối ưu accuracy, mô hình có thể đạt điểm tương đối cao ngay cả khi bỏ sót phần lớn các lớp hiếm nhưng có ý nghĩa an toàn rõ rệt. Vì thế, việc đánh giá bằng \textit{macro-F1} không chỉ là lựa chọn kỹ thuật, mà còn là một quyết định gắn trực tiếp với mục tiêu ứng dụng: không được để các hành vi nguy cơ thấp tần suất bị chìm trong lớp bình thường chiếm ưu thế.

\begin{table}[t]
\caption{So sánh bộ dữ liệu nghiên cứu với một số benchmark liên quan}
\label{tab:dataset-comparison}
\centering
\scriptsize
\begin{tabular}{L{1.2cm}L{0.8cm}L{0.8cm}L{1.1cm}L{2.15cm}}
\toprule
Dataset & Kiểu dữ liệu & Quy mô & Đầu vào & Giá trị đối với bài toán \\
\midrule
DROZY & Drowsy & nhỏ & NIR + sinh lý & Phù hợp cho buồn ngủ sinh lý, hẹp về hành vi cabin tổng quát \\
Drive\&Act & Fine-grain & lớn & RGB, IR, depth, pose & Mạnh về hành vi cabin chi tiết, nhãn rộng, đa view \\
3MDAD & Driving & vừa & RGB + depth, 2 view & Hữu ích cho mất chú ý trong bối cảnh ngày/đêm \\
DBR-DMD-v1.1 & MM clip & 2216 mẫu & body, face, audio, skeleton & Gắn trực tiếp với baseline 5 lớp thống nhất của đề tài \\
\bottomrule
\end{tabular}
\end{table}

Bảng~\ref{tab:dataset-comparison} giúp đặt bộ dữ liệu đang dùng vào đúng bối cảnh. Nếu nhìn theo chiều sâu hành vi, \textit{Drive\&Act} giàu hơn và có phạm vi nhãn rất rộng. Nếu nhìn theo drowsiness sinh lý, \textit{DROZY} lại phù hợp hơn. Tuy nhiên, đối với mục tiêu hiện tại của đề tài, tức là thiết lập một baseline clip-level đa phương thức cho 5 lớp hành vi an toàn tài xế, \textit{DBR-DMD-v1.1 CLEAN\_FULLMM\_5CLASS} là lựa chọn hợp lý hơn vì nhãn, định dạng input và giao thức đánh giá gắn trực tiếp với bài toán đang giải.

\subsection{Hợp đồng đầu vào và tiền xử lý}
Để báo cáo không nhầm lẫn giữa ``artifact đã xác minh'' và ``chi tiết tái dựng hợp lý'', phần này tách rõ hai lớp thông tin. Lớp thứ nhất là \textit{input contract} có thể đọc trực tiếp từ checkpoint và các script đi kèm: hai nhánh hình ảnh \textit{body} và \textit{face} nhận clip cố định 16 frame; nhánh \textit{audio} nhận đặc trưng \textit{log-mel} kích thước $64 \times 40$; nhánh \textit{skeleton} nhận tensor kích thước $16 \times 50$. Lớp thứ hai là cách diễn giải tiền xử lý ở mức phương pháp luận: một video cabin phải được chia thành các cửa sổ thời gian đồng bộ, từ đó sinh ra bốn luồng tín hiệu ăn khớp với nhau trước khi đưa vào mô hình. Báo cáo chỉ khẳng định chắc lớp thứ nhất; lớp thứ hai được dùng để làm rõ pipeline nghiên cứu mà không overclaim những siêu tham số tiền xử lý cuối cùng chưa được đóng gói đầy đủ trong bundle.

\begin{table}[t]
\caption{Hợp đồng đầu vào của mô hình đa phương thức thống nhất}
\label{tab:input-contract}
\centering
\scriptsize
\begin{tabular}{L{0.85cm}L{1.15cm}L{1.2cm}L{0.95cm}L{1.65cm}}
\toprule
Nhánh & Nguồn thô & Biểu diễn đồng bộ & Kích thước & Vai trò chính \\
\midrule
Body & video cabin & clip RGB thân trên & 16 frame & tư thế lái, vai-tay, thao tác lớn \\
Face & video cabin & clip RGB khuôn mặt & 16 frame & mắt, miệng, góc đầu, biểu hiện mặt \\
Audio & sóng âm cabin & log-mel spectrogram & $64 \times 40$ & lời nói và tín hiệu âm thanh bổ trợ \\
Skeleton & landmark theo frame & chuỗi landmark chuẩn hóa & $16 \times 50$ & động học đầu-cổ-vai theo thời gian \\
\bottomrule
\end{tabular}
\end{table}

Ở mức bài toán, một mẫu không nên được hiểu như ``một ảnh kèm vài đặc trưng phụ'', mà như một \textit{clip đa phương thức đồng bộ}. Nếu ký hiệu video đầu vào là $V$, một cửa sổ thời gian bất kỳ được rút gọn thành mẫu
\begin{equation}
\mathcal{X} = \{X^{b}, X^{f}, X^{a}, X^{s}\},
\end{equation}
trong đó $X^{b}$ và $X^{f}$ là hai clip RGB theo thời gian, $X^{a}$ là biểu diễn âm thanh và $X^{s}$ là chuỗi landmark. Cách phát biểu này quan trọng ở hai điểm. Thứ nhất, mọi quyết định cuối cùng đều là quyết định \textit{clip-level}, không phải \textit{frame-level}. Thứ hai, sự đồng bộ thời gian giữa bốn modality là điều kiện bắt buộc để hợp nhất đa nguồn có ý nghĩa.

Về mặt tiền xử lý, luồng dựng mẫu hợp lý cho một video cabin gồm bốn bước: (1) chọn cửa sổ thời gian và lấy các frame đại diện; (2) tạo nhánh \textit{body} từ vùng ảnh cabin chính; (3) tạo nhánh \textit{face} bằng cách định vị và cắt khuôn mặt; (4) trích âm thanh và landmark theo cùng cửa sổ đó. Báo cáo không gắn cứng mọi siêu tham số chưa có trong artifact, nhưng vẫn có thể mô tả đúng bản chất xử lý: \textit{body} và \textit{face} là tín hiệu thị giác theo chuỗi; \textit{skeleton} là tín hiệu hình học theo chuỗi; còn \textit{audio} là tín hiệu một chiều được đưa về miền thời gian-tần số.

Với nhánh âm thanh, nếu $x^{a}$ là waveform của một cửa sổ video, biểu diễn đầu vào được đưa về miền \textit{log-mel} dưới dạng
\begin{equation}
\mathrm{Mel}(f) = 2595 \log_{10}\left(1 + \frac{f}{700}\right),
\end{equation}
\begin{equation}
X^{a} = \log \big(\mathrm{MelSpec}(x^{a}) + \epsilon\big),
\end{equation}
trong đó $\epsilon > 0$ được dùng để ổn định phép lấy log. Hai công thức này không nhằm công bố lại chi tiết toàn bộ script tiền xử lý, mà để làm rõ bản chất của đầu vào \textit{audio}: mô hình không ăn trực tiếp waveform thô, mà ăn một biểu diễn cô đọng đồng thời thông tin thời gian và phổ tần phù hợp với nhận dạng hành vi trong cabin.

Với nhánh skeleton, một clip được xem như chuỗi điểm mốc
\begin{equation}
X^{s} \in \mathbb{R}^{T \times J \times C},
\end{equation}
trong đó $T=16$ là số bước thời gian, còn tích $J \times C = 50$ phản ánh số giá trị landmark sau khi làm phẳng theo tọa độ chuẩn hóa. Điểm mấu chốt ở đây không phải con số 50 tự thân, mà là việc landmark đã được đưa về một chuẩn hình học nhất quán để mô hình học được động học tương đối thay vì bị lệ thuộc quá mạnh vào vị trí tuyệt đối trong khung hình.

\begin{figure*}[t]
\caption{Từ video cabin đến mẫu clip đa phương thức đồng bộ dùng cho suy luận}
\label{fig:preprocessing-pipeline}
\centering
\resizebox{\textwidth}{!}{%
\begin{tikzpicture}[
font=\small,
box/.style={draw, rounded corners=2pt, align=center, minimum height=9mm, inner sep=4pt},
src/.style={box, fill=blue!8, minimum width=23mm},
pre/.style={box, fill=cyan!7, minimum width=33mm},
branch/.style={box, fill=orange!8, minimum width=31mm},
artifact/.style={box, fill=green!8, minimum width=29mm},
arrow/.style={-{Latex[length=2mm]}, thick, draw=blue!55!black},
note/.style={font=\scriptsize, text=black!75}
]
\node[src] (video) at (0,0) {\textbf{Video cabin}\\RGB + audio};
\node[pre] (window) at (3.5,0) {\textbf{Cửa sổ clip}\\đồng bộ theo thời gian};
\node[pre] (sample) at (7.2,0) {\textbf{Lấy mẫu theo clip}\\16 mốc thời gian};

\node[branch] (bodypre) at (10.9,2.1) {\textbf{Nhánh body}\\frame cabin / thân trên};
\node[branch] (facepre) at (10.9,0.7) {\textbf{Nhánh face}\\crop khuôn mặt};
\node[branch] (audiopre) at (10.9,-0.7) {\textbf{Nhánh audio}\\cắt sóng âm đồng bộ};
\node[branch] (skelpre) at (10.9,-2.1) {\textbf{Nhánh skeleton}\\landmark theo frame};

\node[artifact] (bodyx) at (15.0,2.1) {$X^b$\\clip RGB 16 frame};
\node[artifact] (facex) at (15.0,0.7) {$X^f$\\clip RGB 16 frame};
\node[artifact] (audiox) at (15.0,-0.7) {$X^a$\\log-mel $64 \times 40$};
\node[artifact] (skelx) at (15.0,-2.1) {$X^s$\\tensor $16 \times 50$};

\draw[arrow] (video) -- (window);
\draw[arrow] (window) -- (sample);
\draw[arrow] (sample.east) -- ++(0.8,0) |- (bodypre.west);
\draw[arrow] (sample.east) -- ++(0.8,0) |- (facepre.west);
\draw[arrow] (sample.east) -- ++(0.8,0) |- (audiopre.west);
\draw[arrow] (sample.east) -- ++(0.8,0) |- (skelpre.west);
\draw[arrow] (bodypre.east) -- (bodyx.west);
\draw[arrow] (facepre.east) -- (facex.west);
\draw[arrow] (audiopre.east) -- (audiox.west);
\draw[arrow] (skelpre.east) -- (skelx.west);

\node[note, above=1.5mm of bodyx] {thị giác thân người};
\node[note, above=1.5mm of facex] {thị giác khuôn mặt};
\node[note, below=1.5mm of audiox] {thời gian--tần số};
\node[note, below=1.5mm of skelx] {hình học theo chuỗi};
\end{tikzpicture}
}
\end{figure*}

Sơ đồ tiền xử lý ở Hình~1 được vẽ lại theo hướng hệ thống hơn: người đọc không chỉ thấy ``bốn modality đi vào fusion'', mà thấy rõ quá trình từ video thô đến \textit{sample} nghiên cứu. Đây là điểm cần thiết cho một báo cáo học thuật nghiêm túc, vì nó khóa lại đơn vị phân tích thực sự của mô hình. Không phải video toàn cục, cũng không phải frame rời rạc, mà là các \textit{cửa sổ clip đồng bộ} được ánh xạ thành bốn tensor đầu vào có nghĩa.

Ở mức ngữ nghĩa, bốn modality này không cạnh tranh nhau mà bù trừ cho nhau. Nhánh \textit{face} mang tín hiệu rõ nhất cho \textit{eye\_closure} và \textit{yawning}; \textit{body} và \textit{skeleton} mạnh hơn ở các tình huống cần đọc hướng chú ý hoặc thay đổi tư thế vùng thân trên; \textit{audio} thường không phải nguồn quyết định chính nhưng vẫn hữu ích khi tín hiệu thị giác yếu hoặc không ổn định. Chính logic bù trừ này là lý do mô hình được thiết kế theo hướng hợp nhất có điều kiện thay vì coi mọi modality luôn quan trọng như nhau.

\section[Kien Truc De Xuat]{Kiến Trúc Đề Xuất}
Mô hình trung tâm của báo cáo được xem như một \textit{baseline đa phương thức thống nhất ở mức clip}. Đầu vào không phải là ảnh rời rạc, mà là một đơn vị suy luận đồng bộ gồm bốn modality: \textit{body}, \textit{face}, \textit{audio} và \textit{skeleton}. Mỗi modality đi qua một nhánh mã hóa riêng, sau đó được chiếu về cùng không gian biểu diễn để hợp nhất có điều kiện, rồi mới tích lũy ngữ cảnh theo thời gian và suy ra xác suất lớp. Cách phát biểu này khóa lại đúng đối tượng nghiên cứu: không phải bài toán nhận dạng khung hình đơn lẻ, mà là bài toán tích lũy bằng chứng hành vi trên một clip ngắn có ngữ cảnh.

Về mặt phương pháp, kiến trúc phải trả lời bốn câu hỏi liên tiếp: mỗi nguồn tín hiệu được biểu diễn theo dạng nào; các nhánh được đưa về cùng không gian đặc trưng bằng cách nào; cơ chế nào điều tiết mức tin cậy khi một số modality mạnh hơn các modality còn lại; và bằng chứng theo chuỗi được tích lũy ra sao trước khi ra quyết định cuối. Toàn bộ phần Method vì vậy được tổ chức theo đúng bốn tầng đó, thay vì mô tả mô hình như một danh sách backbone đặt cạnh nhau.

\begin{figure*}[t]
\caption{Kiến trúc baseline: mã hóa riêng bốn modality, chiếu về cùng không gian, gated fusion, temporal head và phân lớp 5 hành vi}
\label{fig:model-architecture}
\centering
\resizebox{\textwidth}{!}{%
\begin{tikzpicture}[
font=\small,
box/.style={draw, rounded corners=2pt, align=center, minimum height=8.5mm, inner sep=4pt},
input/.style={box, fill=blue!8, minimum width=23mm},
enc/.style={box, fill=orange!10, minimum width=31mm},
core/.style={box, fill=red!9, minimum width=33mm},
outputbox/.style={box, fill=green!10, minimum width=30mm},
arrow/.style={-{Latex[length=2mm]}, thick, draw=blue!55!black},
stage/.style={font=\scriptsize\bfseries, text=black!75}
]
\node[input] (xb) at (0,2.4) {$X^b$\\body clip};
\node[input] (xf) at (0,0.8) {$X^f$\\face clip};
\node[input] (xa) at (0,-0.8) {$X^a$\\log-mel};
\node[input] (xs) at (0,-2.4) {$X^s$\\skeleton};

\node[enc] (eb) at (4.0,2.4) {\textbf{Body encoder}\\MobileNetV3-Small};
\node[enc] (ef) at (4.0,0.8) {\textbf{Face encoder}\\MobileNetV3-Small};
\node[enc] (ea) at (4.0,-0.8) {\textbf{Audio encoder}\\MLP + GRU};
\node[enc] (es) at (4.0,-2.4) {\textbf{Skeleton encoder}\\MLP + BiGRU};

\node[core] (fusion) at (8.35,0) {\textbf{Gated fusion}\\ước lượng độ tin cậy theo modality};
\node[core] (temporal) at (12.55,0) {\textbf{BiGRU + attention}\\suy luận trình tự theo clip};
\node[outputbox] (head) at (16.6,0) {\textbf{5-class classifier}\\normal / eye / yawn / away / phone};

\draw[arrow] (xb.east) -- (eb.west);
\draw[arrow] (xf.east) -- (ef.west);
\draw[arrow] (xa.east) -- (ea.west);
\draw[arrow] (xs.east) -- (es.west);
\draw[arrow] (eb.east) -- ++(0.6,0) |- (fusion.west);
\draw[arrow] (ef.east) -- ++(0.6,0) |- (fusion.west);
\draw[arrow] (ea.east) -- ++(0.6,0) |- (fusion.west);
\draw[arrow] (es.east) -- ++(0.6,0) |- (fusion.west);
\draw[arrow] (fusion.east) -- (temporal.west);
\draw[arrow] (temporal.east) -- (head.west);

\node[stage] at (0,3.35) {Đầu vào đồng bộ};
\node[stage] at (4.0,3.35) {Mã hóa theo modality};
\node[stage] at (8.35,1.45) {Hợp nhất đa nguồn};
\node[stage] at (12.55,1.45) {Temporal reasoning};
\node[stage] at (16.6,1.45) {Phân loại hành vi};
\end{tikzpicture}
}
\end{figure*}

Hình~\ref{fig:model-architecture} được dựng lại như một sơ đồ suy luận nhiều tầng thay vì một hình ghép backbone. Người đọc nhìn thấy rõ năm khối chức năng nối tiếp nhau: đầu vào đồng bộ, mã hóa theo modality, chiếu và hợp nhất có điều kiện, tích lũy ngữ cảnh theo thời gian, rồi mới đến đầu phân loại. Cách vẽ này khiến sơ đồ trả lời được câu hỏi phương pháp, không chỉ minh họa rằng ``có bốn nhánh đi vào một mô hình''.

\subsection{Nguyên tắc thiết kế kiến trúc}
Kiến trúc được xây theo bốn nguyên tắc. Thứ nhất là \textit{đúng modality}: mỗi nguồn tín hiệu được giữ ở dạng phù hợp với bản chất dữ liệu thay vì ép toàn bộ về một encoder duy nhất. Thứ hai là \textit{projection trước khi hợp nhất}: các embedding riêng phải được đưa về cùng không gian đặc trưng để phép hợp nhất có ý nghĩa về hình học. Thứ ba là \textit{gating theo clip}: mô hình không giả định bốn modality đóng góp ngang nhau trong mọi điều kiện, mà học mức tin cậy theo ngữ cảnh của từng clip. Thứ tư là \textit{quyết định theo chuỗi}: nhãn hành vi chỉ được tạo ra sau khi bằng chứng đã đi qua tầng tích lũy thời gian.

Các nguyên tắc này phù hợp với quy mô và mức độ đóng gói của artifact đang có. Nếu dùng kiến trúc quá nông, mô hình sẽ bỏ sót quan hệ giữa các nguồn tín hiệu. Nếu dùng backbone video quá lớn cho toàn hệ, rủi ro quá khớp và chi phí huấn luyện sẽ tăng nhanh trên tập dữ liệu chỉ hơn hai nghìn mẫu. Thiết kế hiện tại vì vậy có chủ ý giữ ở mức ``đủ sâu để mô hình hóa tương tác'' nhưng vẫn ``đủ gọn để đọc, kiểm chứng và phân tích lỗi''.

\subsection{Chi tiết hơn về fusion và temporal reasoning}
Ký hiệu $E_m(\cdot)$ là encoder của modality $m \in \{b,f,a,s\}$. Sau bước mã hóa riêng, mỗi đầu vào được ánh xạ thành một chuỗi đặc trưng
\begin{equation}
H_m = E_m(X^m) = [h^m_1, h^m_2, \ldots, h^m_T].
\end{equation}
Ở đây cần tách rõ hai vai trò. Tầng encoder chịu trách nhiệm rút đặc trưng riêng theo modality; tầng fusion chịu trách nhiệm làm cho các đặc trưng đó trở nên so sánh được. Hai nhánh thị giác \textit{body} và \textit{face} dùng \textit{MobileNetV3-Small}~\cite{mobilenetv3} để giữ cân bằng giữa năng lực biểu diễn và chi phí tính toán. Nhánh \textit{audio} dùng \textit{MLP + GRU}~\cite{gru} trên biểu diễn \textit{log-mel}, còn nhánh \textit{skeleton} dùng \textit{MLP + BiGRU} để giữ được động học landmark theo chuỗi.

Do bốn nhánh sinh ra embedding dị loại, một bước chiếu về cùng không gian đặc trưng là cần thiết trước khi hợp nhất. Ký hiệu $\phi_m(\cdot)$ là phép chiếu gọn đưa đặc trưng của modality $m$ về không gian có cùng số chiều:
\begin{equation}
\bar{h}^{m}_{t} = \phi_m(h^m_t), \qquad \bar{h}^{m}_{t} \in \mathbb{R}^{d}.
\end{equation}
Nếu bỏ bước này, phép cộng hoặc gating liên-modality chỉ còn là thao tác hình thức vì từng nhánh vẫn nằm ở những thang biểu diễn khác nhau.

Sau khi từng modality đã có embedding riêng, tầng fusion không cộng gộp cứng mà học một hệ số cổng cho từng nhánh:
\begin{equation}
g_m = \sigma(W_m \bar{h}_m + b_m),
\end{equation}
\begin{equation}
\tilde{h}_m = g_m \odot \bar{h}_m,
\end{equation}
trong đó $\bar{h}_m$ là biểu diễn đã được chiếu về cùng không gian đặc trưng, $\sigma(\cdot)$ là hàm sigmoid, và $\odot$ là phép nhân từng phần tử. Có thể hiểu $g_m$ như một phép ước lượng mức nên tin của modality $m$ trong clip đang xét. Khi đó, vector hợp nhất tại bước thời gian $t$ được viết là
\begin{equation}
z_t = \sum_{m \in \{b,f,a,s\}} \tilde{h}^{m}_t.
\end{equation}
Ý nghĩa của tầng này rất quan trọng trong bối cảnh cabin: có clip mà bằng chứng mạnh nằm ở khuôn mặt, có clip lại nằm ở tư thế đầu-thân trên, và cũng có clip mà audio gần như không hữu ích. Gated fusion vì vậy biến ``đa nguồn'' thành ``đa nguồn có điều kiện'', giúp mô hình không bị kéo lệch bởi các modality yếu hoặc nhiễu.

Sau fusion, chuỗi biểu diễn chung $Z = [z_1,\ldots,z_T]$ được đưa vào khối suy luận theo thời gian:
\begin{equation}
H = \mathrm{BiGRU}(Z) = [u_1, u_2, \ldots, u_T].
\end{equation}
BiGRU chịu trách nhiệm tích lũy bằng chứng trước-sau trong clip, từ đó chuyển bài toán từ ``nhiều đặc trưng đồng thời'' thành ``một diễn giải hành vi có ngữ cảnh''. Trên đỉnh chuỗi này, attention pooling chọn lọc những đoạn quan trọng nhất:
\begin{equation}
\alpha_t = \frac{\exp(w^\top u_t)}{\sum_{k=1}^{T} \exp(w^\top u_k)},
\end{equation}
\begin{equation}
v = \sum_{t=1}^{T} \alpha_t u_t.
\end{equation}
Nếu BiGRU trả lời câu hỏi ``chuỗi này diễn tiến như thế nào'', thì attention pooling trả lời câu hỏi ``đoạn nào trong chuỗi đáng được nhấn mạnh nhất khi phân lớp''. Sự kết hợp này phù hợp với bài toán cabin vì bằng chứng hành vi thường không phân bố đều trên toàn clip; nhiều tín hiệu chỉ bộc lộ mạnh trong vài thời điểm ngắn.

\begin{table}[t]
\caption{Vai trò của các tầng trong fusion-temporal head}
\label{tab:fusion-temporal}
\centering
\scriptsize
\begin{tabular}{L{1.2cm}L{1.55cm}L{3.1cm}}
\toprule
Tầng & Đầu ra trung gian & Vai trò chính \\
\midrule
Branch encoder & embedding theo modality & Giữ thông tin đặc thù của body, face, audio và skeleton \\
Gated fusion & biểu diễn hợp nhất & Điều chỉnh trọng số đóng góp giữa các modality \\
BiGRU & chuỗi đặc trưng có ngữ cảnh & Mô hình hóa quan hệ trước-sau trong clip \\
Attention pooling & vector clip-level & Nhấn vào đoạn bằng chứng mạnh nhất trước khi phân lớp \\
\bottomrule
\end{tabular}
\end{table}

Bảng~\ref{tab:fusion-temporal} làm rõ vai trò chức năng của từng tầng, từ đó giúp phần Method không bị rơi vào kiểu mô tả ``kể tên backbone'' mà thiếu logic suy luận.

\subsection{Từ biểu diễn đến quyết định lớp}
Vector clip-level $v$ sau attention pooling được đưa qua đầu phân loại cuối cùng:
\begin{equation}
p = \mathrm{softmax}(W_c v + b_c),
\end{equation}
trong đó $p \in \mathbb{R}^{5}$ là phân phối xác suất trên năm lớp hành vi. Do dữ liệu mất cân bằng rõ rệt giữa các lớp, một mục tiêu huấn luyện hợp lý cho bài toán này là cross-entropy có trọng số:
\begin{equation}
\mathcal{L} = - \sum_{c=1}^{C} w_c y_c \log p_c ,
\end{equation}
với $C=5$ và $w_c$ là hệ số lớp. Artifact huấn luyện của phiên bản \textit{v2} còn cho biết mô hình dùng \textit{label smoothing}; khi đó nhãn cứng $y_c$ được thay bằng
\begin{equation}
y'_c = (1-\varepsilon) y_c + \frac{\varepsilon}{C},
\end{equation}
và hàm mất mát được tính trên $y'_c$. Việc viết rõ mục tiêu huấn luyện theo cách này giúp khóa lại bản chất tối ưu hóa của mô hình: không chỉ học ``đúng lớp'', mà còn học theo cách bớt quá tự tin trên các biên quyết định khó.

Một cách nhìn đúng hơn về toàn bộ pipeline là xem đầu ra cuối cùng như kết quả của một chuỗi lọc bằng chứng nhiều tầng. Ở encoder, mô hình hỏi ``trong từng modality đang có tín hiệu gì''. Ở fusion, mô hình hỏi ``trong clip này nên tin modality nào hơn''. Ở temporal head, mô hình hỏi ``bằng chứng đó có đủ bền và đủ nhất quán để đại diện cho cả clip hay không''. Chỉ sau ba bước này nhãn hành vi mới được chốt. Trật tự suy luận đó cũng là chìa khóa để giải thích phần lớn pattern lỗi ở Section VI và VII.

Về mặt đánh giá, diễn giải trên cũng cho thấy một điểm mạnh của baseline: khi mô hình sai, ta vẫn có thể phỏng đoán tầng cơ chế nào đáng nghi nhất. Nếu một lớp có biểu hiện rõ nhưng vẫn bị nhầm rộng, nghi vấn thường nằm ở tầng biểu diễn. Nếu biểu diễn lớp đã đủ mạnh nhưng xác suất cuối cùng vẫn lẫn giữa vài lớp gần nhau, nghi vấn sẽ chuyển sang fusion hoặc temporal reasoning. Vì artifact hiện có thiên về đầu ra xác suất và confusion pattern, báo cáo giữ kết luận ở mức thận trọng: nút thắt của phiên bản hiện tại nằm gần hơn với tầng hợp nhất--quyết định so với tầng cảm nhận thô.

\section[Thiet Lap Huan Luyen Va Danh Gia]{Thiết Lập Huấn Luyện Và Đánh Giá}
Checkpoint trung tâm của báo cáo là \textit{mmrt\_multimodal\_full\_v2\_t4x2}. Dựa trên artifact đã kiểm chứng, quá trình huấn luyện diễn ra trên môi trường Kaggle với GPU \textit{T4x2}. So với \textit{v1}, phiên bản \textit{v2} giữ nguyên hướng mô hình nhưng chọn cách huấn luyện chặt hơn thông qua \textit{DataParallel}, AMP, batch size lớn hơn, \textit{sqrt inverse} class weighting và \textit{label smoothing}. Theo training scaffold đi kèm, metric checkpoint chính là \textit{validation 5-class macro-F1}; test metric chỉ được công bố khi tập test chứa đủ 5 lớp.

\subsection{Câu hỏi thực nghiệm}
Phần thực nghiệm của báo cáo thực chất xoay quanh bốn câu hỏi. Câu hỏi thứ nhất: một baseline đa phương thức thống nhất có thực sự vượt qua giai đoạn nguyên mẫu hay không. Câu hỏi thứ hai: sau khi đã có baseline chạy ổn định, tinh chỉnh huấn luyện còn tạo thêm ích lợi thực sự hay chỉ là nhiễu thống kê. Câu hỏi thứ ba: sai số của mô hình phân bố đồng đều trên mọi lớp hay tập trung ở một biên quyết định cụ thể. Câu hỏi thứ tư: bundle hiện có đủ loại bằng chứng nào để hỗ trợ các kết luận trên. Việc phát biểu bốn câu hỏi này giúp phần \textit{Experiments} mang cấu trúc phản biện rõ hơn: mỗi bảng và mỗi artifact đều được đặt vào vai trò trả lời một câu hỏi nghiên cứu cụ thể, thay vì chỉ đứng như các con số rời rạc.

\begin{table}[t]
\caption{Ba giai đoạn phát triển mô hình}
\label{tab:main-results}
\centering
\small
\begin{tabular}{p{2.65cm}cc}
\toprule
Giai đoạn mô hình & Val macro-F1 & Accuracy \\
\midrule
Smoke prototype & 0.4085 & -- \\
Unified multimodal baseline (v1) & 0.8014 & -- \\
Training-refined proposed model (v2) & 0.8066 & 0.8054 \\
\bottomrule
\end{tabular}
\end{table}

Bảng~\ref{tab:main-results} được trình bày lại theo tên gọi có nghĩa với người đọc thay vì chỉ giữ ký hiệu nội bộ \textit{smoke}, \textit{v1}, \textit{v2}. Cách đặt tên này giúp hội đồng nhìn rõ tiến trình nghiên cứu: từ một nguyên mẫu kiểm tra luồng (\textit{smoke prototype}), mô hình tiến lên một baseline đa phương thức thống nhất chạy ổn định (\textit{v1}), rồi được tinh chỉnh huấn luyện để tạo ra checkpoint đề xuất cuối cùng (\textit{v2}). Bước nhảy lớn nhất xảy ra ở giai đoạn chuyển từ nguyên mẫu sang baseline hợp nhất, xác nhận rằng giá trị cốt lõi của đề tài nằm ở việc thiết lập được pipeline multimodal thống nhất. Giai đoạn từ \textit{v1} sang \textit{v2} chỉ tăng nhẹ 0.0052 macro-F1, cho thấy sau khi lõi mô hình đã thành hình, dư địa cải thiện còn lại nằm nhiều hơn ở biên lớp khó và độ bao phủ thực nghiệm hơn là ở các chỉnh sửa huấn luyện cục bộ.

\begin{table}[t]
\caption{Baseline tham chiếu và checkpoint đề xuất cuối}
\label{tab:baseline-vs-proposed}
\centering
\scriptsize
\begin{tabular}{L{1.75cm}L{2.15cm}L{0.95cm}L{0.8cm}}
\toprule
Phiên bản & Khác biệt chính & Macro-F1 & $\Delta$ \\
\midrule
Baseline hợp nhất (\textit{v1}) & Baseline clip-level đầu tiên chạy ổn định; có log epoch. & 0.8014 & -- \\
Mô hình đề xuất (\textit{v2}) & Giữ lõi kiến trúc, siết huấn luyện và cân bằng lớp. & 0.8066 & +0.0052 \\
\bottomrule
\end{tabular}
\end{table}

Bảng~\ref{tab:baseline-vs-proposed} làm rõ một điểm rất quan trọng cho phần bảo vệ: mô hình đề xuất không phải là một kiến trúc hoàn toàn mới tách rời khỏi baseline, mà là một bước tinh chỉnh có kiểm soát trên một lõi đã chạy được. Điều này khiến đóng góp của đề tài trở nên thuyết phục hơn. Thay vì trình bày như một chuỗi thử nghiệm rời rạc, báo cáo cho thấy một tiến trình nghiên cứu có logic: xác lập baseline, nhận diện giới hạn, rồi nâng chất lượng theo đúng điểm yếu quan sát được.

\subsection{Khả năng tái lập}
Báo cáo phân biệt rõ hai mức độ. Ở mức thứ nhất, có thể tái lập hợp đồng đầu vào, logic chọn checkpoint theo \textit{macro-F1}, các thay đổi huấn luyện chính của \textit{v2}, cùng như artifact dự đoán trên validation. Ở mức thứ hai, một số siêu tham số chi tiết như learning rate cuối, batch size chính xác, số epoch tốt nhất hay chi phí huấn luyện theo epoch chưa được đóng gói đủ trong bundle \textit{full\_v2}. Vì vậy, báo cáo chỉ claim những gì đã xác minh trực tiếp, thay vì điền phần thiếu bằng suy đoán.

\subsection{Giao thức đánh giá}
Do dữ liệu mất cân bằng, \textit{macro-F1} được dùng làm mục tiêu chính để so sánh các phiên bản. Theo scaffold huấn luyện, báo cáo chuẩn cho phép theo dõi thêm \textit{per-class recall}, confusion matrix và balanced accuracy. Trong bundle hiện có, \textit{v1} lưu \textit{history.json} nên có thể quan sát quá trình tiến triển theo epoch; ngược lại, \textit{v2} chỉ lưu checkpoint cuối cùng cùng các bảng dự đoán. Khác biệt này giúp làm rõ thứ bậc bằng chứng: báo cáo có thể khẳng định khá chắc về kết quả validation và pattern lỗi, nhưng chưa đủ bằng chứng để mở rộng thành claim mạnh về test-set generalization, ablation từng modality, hay benchmark realtime.

Ở cấp độ phương pháp luận, việc ưu tiên \textit{macro-F1} thay vì accuracy có ý nghĩa đặc biệt quan trọng với đề tài này. Một hệ thống an toàn tài xế không thể được xem là tốt nếu nó đạt accuracy nhờ làm đúng lớp \textit{normal\_driving} trong khi bỏ sót các lớp hiếm nhưng nguy cơ cao. Vì vậy, báo cáo cố ý xem đây là một bài toán \textit{đánh đổi công bằng giữa các lớp hành vi} hơn là bài toán tối đa hóa tỷ lệ đúng toàn cục. Cách đặt chuẩn đánh giá này cũng làm cho việc phân tích confusion matrix trở thành bước bắt buộc, chứ không phải phần minh họa thêm.

\section[Thuc Nghiem Va Phan Tich Ket Qua]{Thực Nghiệm Và Phân Tích Kết Quả}
\subsection{Cách đọc phần thực nghiệm}
Phần thực nghiệm được viết theo logic ``câu hỏi nghiên cứu $\rightarrow$ bằng chứng $\rightarrow$ kết luận''. Thay vì liệt kê chỉ số theo thứ tự xuất hiện của artifact, báo cáo đi từ câu hỏi tổng quát nhất đến câu hỏi hẹp nhất: mô hình có tốt hơn nguyên mẫu không, phần cải thiện mới nhất có ý nghĩa gì, sai số tập trung ở đâu, và những gì trong bundle đủ mạnh để được xem là bằng chứng. Cách tổ chức này giúp hội đồng nhìn thấy rõ luận điểm trung tâm của đề tài, thay vì chỉ thấy nhiều bảng số nằm cạnh nhau.

\subsection{RQ1: Baseline thống nhất có thực sự vượt qua giai đoạn nguyên mẫu?}
Bằng chứng trực tiếp cho RQ1 nằm ở Bảng~\ref{tab:main-results}. Từ một \textit{smoke prototype} chỉ đạt macro-F1 = 0.4085, pipeline tiến lên mức 0.8014 khi chuyển sang baseline đa phương thức thống nhất \textit{v1}. Đây là bước nhảy lớn nhất trong toàn bộ quá trình phát triển, và vì thế mang ý nghĩa học thuật cao hơn cả bước cải thiện từ \textit{v1} sang \textit{v2}. Nó xác nhận rằng đóng góp cốt lõi của đề tài không nằm ở một mẹo huấn luyện cục bộ, mà nằm ở việc thiết lập được một lõi suy luận clip-level đa phương thức thống nhất.

Nói cách khác, RQ1 được trả lời theo hướng dứt khoát: có, baseline thống nhất đã vượt qua hẳn giai đoạn nguyên mẫu. Nếu không có bước chuyển này, mọi phân tích sâu hơn về fusion, temporal reasoning hay pattern lỗi đều thiếu nền tảng. Giá trị đầu tiên cần chốt với hội đồng vì vậy là sự tồn tại của một baseline nghiên cứu chạy ổn định, chứ không chỉ là một workflow demo ghép nhiều module.

\subsection{RQ2: Việc tinh chỉnh từ v1 sang v2 có tạo ra cải thiện đáng tin?}
Đối với RQ2, bằng chứng nằm ở hai nơi: Bảng~\ref{tab:main-results} và Bảng~\ref{tab:baseline-vs-proposed}. Chênh lệch tuyệt đối từ 0.8014 lên 0.8066 chỉ vào khoảng +0.0052 macro-F1. Nếu đọc hời hợt, người xem có thể cho rằng cải thiện này nhỏ. Tuy nhiên, trong bối cảnh baseline đã ở vùng hiệu năng tương đối cao, ý nghĩa của nó không nằm ở độ lớn tuyệt đối mà ở chỗ hướng cải thiện xuất hiện sau khi kiến trúc cốt lõi đã ổn định.

Artifact \textit{v1} còn cho thấy mô hình ban đầu khá nhạy với quá trình huấn luyện: macro-F1 tăng lên ở epoch 3 nhưng giảm mạnh ở epoch 4. Điều này khiến \textit{v2} nên được đọc như một bước ``siết protocol huấn luyện'' thay vì ``đổi mô hình''. Kết luận đúng cho RQ2 vì thế là: có cải thiện, nhưng đó là cải thiện ở vùng quyết định khó và ở độ ổn định huấn luyện, chứ không phải một bước nhảy kiến trúc mới. Cách diễn giải này trung thực hơn nhiều so với việc cố phóng đại một mức tăng nhỏ thành một đột phá.

\subsection{RQ3: Sai số nằm rải đều hay tập trung ở một biên quyết định?}
RQ3 được trả lời chủ yếu bằng Hình~\ref{fig:confusion-matrix} và Bảng~\ref{tab:error-summary}. Cả hai chỉ ra cùng một thông điệp: sai số không phân bố đều trên toàn bộ không gian lớp. Ba lớp \textit{phone\_use}, \textit{eye\_closure} và \textit{yawning} có tín hiệu nhận diện tương đối sắc, trong khi điểm nghẽn gần như tập trung vào biên \textit{normal\_driving}/\textit{looking\_away}. Cụ thể, 67 trên 83 mẫu \textit{looking\_away} bị kéo về \textit{normal\_driving}.

Điều này có ý nghĩa học thuật rõ ràng. Mô hình hiện tại không thất bại theo kiểu ``yếu đều'', mà thất bại ở một vùng ngữ nghĩa hẹp nhưng rất khó. Với bài toán driver monitoring, đây là thông tin tốt hơn nhiều so với một confusion pattern trải khắp các lớp, vì nó cho thấy hướng cải thiện tiếp theo có thể nhắm đúng vào temporal decision boundary thay vì phải thay toàn bộ mô hình. Nói gọn hơn, RQ3 cho thấy nút thắt của bài toán hiện không còn ở khả năng phát hiện tín hiệu nổi bật, mà ở khả năng tách hai trạng thái gần nhau về ngữ nghĩa và gần nhau về hình thái cục bộ.

\subsection{RQ4: Bundle hiện có chứng minh được gì, và chưa chứng minh được gì?}
RQ4 là câu hỏi về kỷ luật bằng chứng. Bảng~\ref{tab:evidence-scope} cho thấy mỗi loại artifact đóng một vai trò khác nhau: \textit{history.json} của \textit{v1} dùng để đọc động học học tập; \textit{pair predictions} của \textit{v2} dùng để xác định vùng nhầm lẫn; \textit{full probabilities} dùng để nhìn sâu vào độ mờ của quyết định; còn confusion matrix dùng để chốt error pattern ở mức định tính. Nhờ vậy, báo cáo có thể khẳng định khá chắc hiệu năng validation và pattern lỗi, nhưng chưa được phép tuyên bố mạnh về test-set generalization, benchmark realtime hay ablation định lượng của từng modality.

Đây là một điểm cần nói rõ trước hội đồng: bundle hiện có đủ mạnh để bảo vệ một baseline nghiên cứu có kiểm chứng, nhưng chưa đủ dày để bảo vệ các claim kiểu ``audio đóng góp bao nhiêu phần trăm'' hoặc ``runtime app chứng minh xong chất lượng mô hình''. Chính việc tự giới hạn phạm vi kết luận này mới làm phần thực nghiệm đáng tin hơn, không phải yếu đi.

Điều đáng chú ý ở đây là cấu trúc khó của bài toán không phân bố đều giữa các lớp ``nguy hiểm'' và ``không nguy hiểm''. Thay vào đó, mô hình xử lý tốt các lớp có bằng chứng hình thái rõ ngay cả khi chúng là lớp rủi ro, còn vấp ở cặp lớp có biên ngữ nghĩa mờ dù một trong hai là trạng thái bình thường. Phát hiện này giúp báo cáo tránh được một diễn giải quá đơn giản kiểu ``mô hình đa phương thức chỉ giỏi bắt hành vi nổi bật''. Chính xác hơn, mô hình đang cho thấy rằng đa phương thức đặc biệt hữu ích khi có dấu hiệu cục bộ mạnh; còn với các lệch hướng tinh vi, bài toán quyết định theo thời gian vẫn là nút thắt khó nhất.

\subsection{Artifact phân tích sau huấn luyện}
Bundle \textit{v2} lưu hai loại artifact có giá trị phân tích cao: bảng \textit{pair predictions} và bảng \textit{full probabilities}. Bảng dự đoán cho phép xem trực tiếp mẫu nào bị nhầm và nhầm sang lớp nào; bảng xác suất đầy đủ lại cho phép nhìn sâu hơn vào cấu trúc quyết định của mô hình. Nhiều mẫu \textit{normal\_driving} trong \textit{v2} chỉ có xác suất lớp đúng quanh 0.44--0.56, trong khi xác suất \textit{yawning} hoặc \textit{looking\_away} vẫn ở mức 0.22--0.35. Đây là một quan sát quan trọng, vì nó cho thấy sai số chính không đến từ việc representation hoàn toàn thất bại, mà từ những vùng quyết định mờ nơi bằng chứng giữa nhiều lớp gần nhau vẫn chưa được tách đủ rõ. Nói cách khác, mô hình không ``mù'' với hành vi, mà đang thiếu độ sắc ở tầng hợp nhất-thời gian-quyết định.

\subsection{Những gì có thể và chưa thể kết luận}
Từ tập artifact hiện tại, báo cáo có thể kết luận khá chắc bốn điểm: (1) baseline đa phương thức thống nhất hoạt động tốt hơn rõ rệt so với nguyên mẫu; (2) checkpoint \textit{v2} là phiên bản mạnh nhất trong bundle được xác minh; (3) sai số không trải đều mà tập trung ở cặp \textit{normal\_driving}/\textit{looking\_away}; và (4) ba lớp \textit{phone\_use}, \textit{eye\_closure}, \textit{yawning} có tín hiệu ổn định hơn nhiều trong validation. Ngược lại, báo cáo chưa nên kết luận mạnh về đóng góp định lượng của từng modality, chưa nên tuyên bố tổng quát hóa ra ngoài điều kiện validation hiện có, và chưa nên đọc runtime ứng dụng như benchmark cuối cùng của mô hình. Chính việc tách bạch ranh giới này làm phần thực nghiệm đáng tin hơn.

\begin{table}[t]
\caption{Bằng chứng thực nghiệm và phạm vi kết luận}
\label{tab:evidence-scope}
\centering
\scriptsize
\begin{tabular}{L{1.7cm}L{2.45cm}L{1.95cm}}
\toprule
Nguồn & Quan sát chính & Kết luận rút ra \\
\midrule
\textit{history.json} của v1 & Diễn tiến epoch, macro-F1, recall từng lớp & Đọc xu hướng học và độ ổn định baseline \\
Pair predictions của v2 & Mẫu nhầm và lớp bị nhầm sang & Xác định vùng lỗi chi phối \\
Full probabilities của v2 & Phân bố xác suất trên đủ 5 lớp & Phân tích độ tin cậy và vùng quyết định mờ \\
Confusion matrix val & Recall theo lớp và cặp nhầm chính & Chốt error pattern ở mức định tính \\
\bottomrule
\end{tabular}
\end{table}

Bảng~\ref{tab:evidence-scope} làm rõ một nguyên tắc xuyên suốt của báo cáo: mỗi claim đều phải bám vào một loại bằng chứng cụ thể. Cách trình bày này giúp phần thực nghiệm thuyết phục hơn, vì người chấm nhìn được ngay ``mình có gì trong tay'' và ``mình dám kết luận đến đâu''.

\section[Phan Tich Loi]{Phân Tích Lỗi}
Mặc dù kết quả tổng thể khá tốt, hiệu năng không đồng đều trên mọi lớp. Ba lớp \textit{phone\_use}, \textit{eye\_closure} và \textit{yawning} có xu hướng ổn định hơn do mang tín hiệu hình thái rõ. Ngược lại, cặp \textit{normal\_driving} và \textit{looking\_away} tạo ra điểm nghẽn chính của bài toán. Ở \textit{v1}, artifact cho thấy mô hình có xu hướng đẩy một phần mẫu \textit{normal\_driving} sang \textit{looking\_away}; sang \textit{v2}, xu hướng này được giảm mạnh ở nhóm đó nhưng không biến mất hoàn toàn.

\begin{figure}[t]
\caption{Confusion matrix của v2 trên validation; lỗi chính nằm ở biên \textit{looking\_away}/\textit{normal\_driving}}
\label{fig:confusion-matrix}
\centering
\begin{tikzpicture}[x=0.62cm,y=0.62cm]
\node[rotate=90] at (-1.65,2.0) {\scriptsize Nhãn thật};
\node at (2.0,5.95) {\scriptsize Nhãn dự đoán};
\node at (0.5,5.3) {\scriptsize N};
\node at (1.5,5.3) {\scriptsize E};
\node at (2.5,5.3) {\scriptsize Y};
\node at (3.5,5.3) {\scriptsize L};
\node at (4.5,5.3) {\scriptsize P};
\node at (-0.45,4.5) {\scriptsize N};
\node at (-0.45,3.5) {\scriptsize E};
\node at (-0.45,2.5) {\scriptsize Y};
\node at (-0.45,1.5) {\scriptsize L};
\node at (-0.45,0.5) {\scriptsize P};
\foreach \x/\y/\v in {
0/0/176,1/0/0,2/0/2,3/0/3,4/0/0,
0/1/0,1/1/23,2/1/0,3/1/0,4/1/0,
0/2/0,1/2/0,2/2/8,3/2/0,4/2/0,
0/3/67,1/3/0,2/3/0,3/3/16,4/3/0,
0/4/0,1/4/0,2/4/0,3/4/0,4/4/75} {
  \pgfmathsetmacro{\shade}{min(100,int(\v/176*100))}
  \fill[blue!\shade] (\x,4-\y) rectangle ++(1,1);
  \draw[black!35] (\x,4-\y) rectangle ++(1,1);
  \node at (\x+0.5,4.5-\y) {\scriptsize \v};
}
\draw[red, very thick] (0,1) rectangle ++(1,1);
\node[anchor=west, align=left] at (5.25,1.5) {\scriptsize L\,$\rightarrow$\,N là lỗi lớn nhất\\\scriptsize 67 / 83 mẫu \textit{looking\_away}};
\end{tikzpicture}

\vspace{0.2em}
{\scriptsize Chú giải lớp: N = \textit{normal\_driving}, E = \textit{eye\_closure}, Y = \textit{yawning}, L = \textit{looking\_away}, P = \textit{phone\_use}.}
\end{figure}

Hình~\ref{fig:confusion-matrix} cho thấy lỗi chi phối của mô hình không nằm ở các lớp có tín hiệu hình thái rõ, mà tập trung gần như toàn bộ ở lớp \textit{looking\_away}. Cụ thể, 67 trên 83 mẫu \textit{looking\_away} bị dự đoán thành \textit{normal\_driving}, cho thấy đây mới là ranh giới quyết định khó nhất của baseline hiện tại.

\begin{table}[t]
\caption{Tóm tắt vùng lỗi chính trên validation}
\label{tab:error-summary}
\centering
\scriptsize
\begin{tabular}{p{1.45cm}rrp{2.15cm}}
\toprule
Lớp & Support & Recall & Lỗi chi phối \\
\midrule
Normal & 181 & 0.972 & Rải rác sang Y hoặc L \\
Eye closure & 23 & 1.000 & Không đáng kể \\
Yawning & 8 & 1.000 & Không đáng kể \\
Looking away & 83 & 0.193 & 67 mẫu sang N \\
Phone use & 75 & 1.000 & Không đáng kể \\
\bottomrule
\end{tabular}

\vspace{0.15em}
{\scriptsize Ký hiệu N, Y, L được giữ như trong Hình~\ref{fig:confusion-matrix}.}
\end{table}

Bảng~\ref{tab:error-summary} nén lại thông điệp của confusion matrix dưới dạng dễ chấm hơn: bài toán hiện không thất bại đồng đều, mà thất bại tập trung ở đúng một vùng ngữ nghĩa hẹp. Với hội đồng, đây là một dấu hiệu tốt hơn nhiều so với trường hợp mô hình yếu trên toàn bộ các lớp, vì nó cho thấy hướng cải thiện tiếp theo đã khá rõ.

Mẫu nhầm lẫn này có thể giải thích trên ba mức. Ở mức dữ liệu, ranh giới giữa hai hành vi có thể mơ hồ khi góc quay, chất lượng khung hình hoặc thời lượng lệch nhìn không đủ lớn. Ở mức biểu diễn, tín hiệu cần cho cặp lớp này tinh vi hơn so với \textit{phone\_use} hay \textit{eye\_closure}; khác với điện thoại hoặc mắt nhắm, \textit{looking\_away} có thể chỉ biểu hiện qua thay đổi nhỏ của hướng đầu, ánh mắt hoặc tư thế thân trên trong một khoảng rất ngắn. Ở mức ra quyết định, clip ngắn 16 frame có thể chưa đủ để làm bằng chứng theo thời gian ổn định trong một số tình huống biên, đặc biệt khi phần lớn frame vẫn gần với trạng thái lái bình thường.

Về mặt nghiên cứu, phát hiện này có hai ý nghĩa. Thứ nhất, nó chỉ ra rằng việc nâng tiếp hiệu năng nhiều khả năng không nằm ở việc tăng độ phức tạp mô hình một cách mù quáng, mà nằm ở cách mô hình hóa bằng chứng thời gian cho nhóm lớp gần nhau. Thứ hai, nó cung cấp hướng ưu tiên rõ ràng cho các thực nghiệm tiếp theo: \textit{temporal decision smoothing}, calibration, persistence rule, hoặc các loss/weighting nhấn mạnh cặp lớp \textit{normal\_driving}/\textit{looking\_away}. Tuy nhiên, đây vẫn là đề xuất rút ra từ pattern lỗi, không phải claim đã được chứng minh bằng thực nghiệm bổ sung trong bundle hiện có.

Một giới hạn quan trọng cần nói rõ là bundle chưa cung cấp ablation hoàn chỉnh theo từng modality cho chính \textit{mmrt\_multimodal\_full\_v2\_t4x2}. Do đó, báo cáo không kết luận định lượng rằng bỏ \textit{audio} hay \textit{skeleton} sẽ làm giảm bao nhiêu \textit{macro-F1}; phần đánh giá ở đây chỉ dám khẳng định vai trò thiết kế của từng modality và các pattern lỗi quan sát được.

\section[Tang Ung Dung Hoa MDDAD]{Minh Họa Ứng Dụng Hóa Bằng MDDAD}
\textit{MDDAD} chỉ được đặt trong báo cáo như một \textit{application layer}, không phải đối tượng đánh giá học thuật trung tâm. Trục chính của nghiên cứu vẫn là baseline \textit{mmrt\_multimodal\_full\_v2\_t4x2}; vì vậy mọi kết luận về chất lượng nhận diện 5 lớp phải bám vào artifact của baseline này. Vai trò đúng của \textit{MDDAD} là chứng minh rằng một lõi nghiên cứu đã được kiểm chứng có thể được tổ chức thành workflow phục vụ upload analysis, monitoring và reporting.

Trong trạng thái hiện tại, hệ thống duy trì ba nhóm chức năng: phân tích video tải lên, giám sát webcam gần thời gian thực, và tổng hợp đánh giá rủi ro theo phiên. Route \texttt{POST /sessions/<id>/multimodal-analysis} cho thấy nhánh upload đã được tổ chức quanh suy luận theo clip và theo cửa sổ thời gian. Tuy nhiên, điều đó không đồng nghĩa ứng dụng đã triển khai đầy đủ baseline \textit{mmrt}; nó chỉ cho thấy lớp ứng dụng đang cố bảo toàn tư duy suy luận theo ngữ cảnh thời gian khi chuyển từ nghiên cứu sang sản phẩm demo.

\begin{table}[t]
\caption{Workload suy luận đã xác minh của application layer}
\label{tab:runtime-config}
\centering
\scriptsize
\begin{tabular}{L{1.55cm}L{0.75cm}L{3.45cm}}
\toprule
Tham số & Giá trị & Ý nghĩa triển khai \\
\midrule
Clip window & 4.0s & Mỗi quyết định dựa trên một cửa sổ dài 4 giây. \\
Clip stride & 2.0s & Tạo xấp xỉ 30 quyết định clip trong một phút video. \\
Frames/clip & 16 & Tương đương khoảng 4 frame mẫu mỗi giây cho nhánh hành vi. \\
Realtime FPS & 12.0 & Vòng lặp webcam ưu tiên độ ổn định hiển thị. \\
Preview FPS & 3.0 & Luồng preview giảm tải so với camera đầy đủ. \\
Upload FPS & 5.0 & Nhánh upload lấy mẫu thưa hơn để giữ chi phí xử lý hợp lý. \\
\bottomrule
\end{tabular}
\end{table}

Bảng~\ref{tab:runtime-config} chỉ báo cáo \textit{workload suy luận đã xác minh}, không tuyên bố benchmark latency. Nói cách khác, đây là các đại lượng chi phí triển khai có thể kiểm chứng trực tiếp từ cấu hình và pipeline; còn thời gian suy luận tính bằng mili giây hay FPS chuẩn benchmark vẫn cần một thí nghiệm riêng với phần cứng, batch size và chế độ chạy được khóa chặt.

Điểm cần phân biệt dứt khoát là nhánh giám sát hành vi thời gian thực trong \textit{MDDAD} hiện vẫn chủ yếu dựa trên MediaPipe, EAR/MAR, head-pose heuristic và các temporal rule; tài liệu hệ thống cũng ghi rõ \textit{Built-in MediaPipe/OpenCV rules} và \textit{Chưa có model action recognition train riêng}. Nhánh học sâu rõ nhất trong app hiện hành là phát hiện điện thoại bằng detector, trong đó RT-DETR-L serious fine-tuned là model chính, bên cạnh các phương án RT-DETR compare và YOLO26n fallback. Do đó, \textit{MDDAD} không được dùng để chứng minh chất lượng của bài toán phân loại đa phương thức 5 lớp; nó chỉ minh họa lộ trình ứng dụng hóa sau nghiên cứu.

\subsection{Phạm vi liên hệ với baseline nghiên cứu}
Điểm giao nhau quan trọng nhất giữa \textit{MDDAD} và baseline \textit{mmrt} nằm ở tư duy suy luận theo clip và theo thời gian, chứ không nằm ở việc hai bên đã dùng chung toàn bộ kiến trúc nội bộ. Baseline nghiên cứu giải quyết trực tiếp bài toán 5 lớp trên dữ liệu đa phương thức đồng bộ. Ngược lại, app hiện hành ưu tiên demo ổn định nên dùng pipeline lai ghép, trong đó phần hành vi vẫn thiên về tín hiệu landmark và temporal rule, còn detector điện thoại là một nhánh bổ sung. Việc tách hai tầng theo đúng bản chất này giúp báo cáo không lẫn giữa đóng góp học thuật và đóng góp triển khai.

\subsection{Đối chiếu baseline nghiên cứu và runtime ứng dụng}
Một điểm dễ gây hiểu nhầm là sự khác nhau giữa baseline nghiên cứu \textit{mmrt\_multimodal\_full\_v2\_t4x2} và runtime hiện hành của \textit{MDDAD}. Để tránh nhập nhằng, Bảng~\ref{tab:research-vs-app} đối chiếu hai tầng này theo nguồn dữ liệu, đầu vào, cơ chế quyết định và loại bằng chứng đánh giá.

\begin{table}[t]
\caption{Baseline nghiên cứu và runtime hiện hành của MDDAD}
\label{tab:research-vs-app}
\centering
\small
\begin{tabular}{L{1.8cm}L{2.6cm}L{2.6cm}}
\toprule
Tiêu chí & Baseline nghiên cứu & Runtime app hiện hành \\
\midrule
Nguồn chính & DBR-DMD-v1.1 5 lớp & Webcam/upload cục bộ \\
Đầu vào & body, face, audio, skeleton & frame video, landmark, detector điện thoại \\
Đơn vị suy luận & clip 16 frame đồng bộ & frame + luật thời gian theo cửa sổ \\
Quyết định & gated fusion + BiGRU attention & heuristic + detector + risk engine \\
Kết quả xác minh & val macro-F1 = 0.8066 & workflow chạy được, có score và report \\
Điểm mạnh & chuẩn nghiên cứu, có artifact xác suất & ổn định demo, có dashboard và history \\
Giới hạn & thiếu test metric đóng gói & chưa là unified multimodal runtime \\
\bottomrule
\end{tabular}
\end{table}

Bảng~\ref{tab:research-vs-app} khóa lại ranh giới diễn giải của báo cáo. Cột trái là phần có giá trị nghiên cứu và được hỗ trợ bởi artifact thực nghiệm; cột phải là tầng ứng dụng hóa và tuyệt đối không được dùng để suy diễn ngược về chất lượng của baseline. Điểm này cần giữ thật sắc, vì nếu trộn hai tầng vào nhau, phần bảo vệ rất dễ rơi vào thế ``lấy demo app để chứng minh mô hình''.

\subsection{Luồng vận hành của application layer}
Theo tài liệu kiến trúc, cả workflow upload và workflow realtime đều đi qua cùng một vòng đời mức cao: nhận đầu vào, mở session, sinh tín hiệu theo thời gian, tổng hợp event, cập nhật score và lưu kết quả để hiển thị lại trên dashboard hoặc report. Mô tả này đủ để cho thấy \textit{MDDAD} có cấu trúc ứng dụng rõ ràng, nhưng không đẩy báo cáo sang hướng phân tích chi tiết một pipeline khác với baseline nghiên cứu. Vì mục tiêu chính của đề tài là mô hình đa phương thức, phần app chỉ nên dừng ở mức chứng minh tính khả dụng của hướng chuyển giao.

\section[Kiem Thu He Thong Va Tieu Chi Chap Nhan]{Kiểm Thử Hệ Thống Và Tiêu Chí Chấp Nhận}
Section này chỉ áp dụng cho tầng ứng dụng \textit{MDDAD}. Nó không mở rộng, thay thế hay hiệu chỉnh kết luận của phần thực nghiệm mô hình. Bằng chứng cốt lõi cho baseline nghiên cứu vẫn chỉ gồm \textit{validation macro-F1 = 0.8066}, confusion matrix validation và các artifact dự đoán. Vì vậy, toàn bộ nội dung ở đây chỉ trả lời một câu hỏi mang tính triển khai: application layer có chạy trọn workflow và sinh được đầu ra dùng được cho demo hay không.

Theo tài liệu \textit{evaluation plan}, phần nghiệm thu hệ thống nên tập trung vào bốn nhóm kiểm tra: workflow chức năng đầu-cuối, khả năng giải thích của timeline và report, các chỉ số vận hành đơn giản như score hay event count, và các failure case có kiểm soát. Đây là tiêu chí chấp nhận của sản phẩm demo, không phải benchmark học thuật của mô hình \textit{mmrt}. Vì thế, các thành phần như detector fallback, CSV export hay camera error chỉ được dùng để đánh giá độ hoàn chỉnh của lớp ứng dụng, không được dùng để suy diễn ngược về chất lượng của baseline nghiên cứu.

\subsection{Protocol đánh giá theo kiểu đồ án tốt nghiệp}
Từ requirement specification và evaluation plan, protocol đánh giá phù hợp cho đồ án này nên chia thành bốn lớp: đánh giá chức năng đầu-cuối, đánh giá định tính về khả năng giải thích, đánh giá định lượng đơn giản trên phiên chạy thật, và đánh giá failure case có kiểm soát. Cách chia này phù hợp cho phần \textit{ứng dụng hóa}, miễn là được đặt tách biệt khỏi phần thực nghiệm của baseline nghiên cứu.

Lớp thứ nhất là \textit{functional end-to-end evaluation}. Theo evaluation plan, cần kiểm tra trọn vẹn các workflow của app như login/logout, tạo driver, uploaded video creation, MediaPipe analysis, phone-detector workflow, realtime monitoring, snapshot preview, dashboard, history, reports, evaluation page và CSV export. Với mỗi workflow, cách ghi nhận nên theo mẫu \textit{Expected Result / Actual Result / Pass-Fail / Notes}. Đây là kiểu nghiệm thu phù hợp cho bảo vệ tốt nghiệp của \textit{MDDAD}, nhưng nó không thay thế cho đánh giá nhận diện 5 lớp của baseline nghiên cứu.

Lớp thứ hai là \textit{qualitative explainability evaluation}. Vì project tự định vị là local-first và explainable, đánh giá không thể bỏ qua câu hỏi liệu người xem có hiểu được timeline, evidence snapshot, risk explanation và report hay không. Các câu hỏi trong evaluation plan như ``event label có dễ hiểu không'', ``timeline có khớp hành vi nhìn thấy không'', hay ``reviewer có theo dõi được cách risk biến thành safety score và report không'' là các tiêu chí rất sát tinh thần đồ án. Chúng đặc biệt quan trọng với \textit{MDDAD}, bởi một hệ thống minh bạch nhưng khó diễn giải trên giao diện thì vẫn chưa hoàn thành mục tiêu ứng dụng hóa.

Lớp thứ ba là \textit{simple quantitative evaluation}. Trong phạm vi này, các chỉ số thực tế nên được trích từ evaluation page và CSV export, bao gồm total sessions, safety score, risk score, alert level, event count, event count by type, processed frame count, face detected frame count, face detection rate, approximate processing FPS và phone usage count. Chúng chỉ mang ý nghĩa vận hành của app. Chẳng hạn, face detection rate được dùng để giải thích false negative của MediaPipe, còn event count by type được dùng để kiểm tra video chuẩn có tạo ra loại event kỳ vọng hay không.

\begin{table}[t]
\caption{Protocol nghiệm thu lớp ứng dụng}
\label{tab:evaluation-protocol}
\centering
\scriptsize
\begin{tabular}{L{1.45cm}L{2.65cm}L{2.15cm}}
\toprule
Nhóm & Trọng tâm kiểm tra & Bằng chứng \\
\midrule
Chức năng & đăng nhập, driver, upload, realtime, báo cáo, export & checklist đạt/không đạt \\
Định tính & timeline, giải thích, evidence, độ rõ báo cáo & ghi chú của người đánh giá \\
Định lượng & score, số event, face rate, FPS, phone count & trang evaluation và CSV \\
Chịu lỗi & lỗi camera, file sai, thiếu detector, không thấy mặt & hành vi fallback \\
\bottomrule
\end{tabular}
\end{table}

Lớp thứ tư là \textit{controlled failure-case evaluation}. Theo requirement và evaluation plan, các ca lỗi tối thiểu cần được ghi nhận gồm camera unavailable, missing detector model, missing uploaded video file, no face detected, poor lighting và invalid upload file. Giá trị của lớp đánh giá này nằm ở chỗ nó chứng minh application layer không chỉ chạy khi mọi điều kiện thuận lợi, mà còn xuống cấp một cách có kiểm soát.

\subsection{Tiêu chí chấp nhận cuối kỳ}
Dựa trên các yêu cầu chức năng và phi chức năng, có thể chốt bộ tiêu chí chấp nhận cuối kỳ của application layer ở ba mức. Ở mức \textit{bắt buộc}, hệ thống phải chạy local, đăng nhập được, hỗ trợ nhiều driver, phân tích được webcam hoặc video upload, sinh được event, score, history và report, đồng thời không phụ thuộc dịch vụ trả phí. Ở mức \textit{nên có}, hệ thống nên hiển thị live update, fatigue trend forecast và filtering hợp lý. Ở mức \textit{mở rộng}, evidence snapshot, report export, calibration và optional detector backend sẽ làm tăng chất lượng bảo vệ nhưng không phải điều kiện để đánh giá baseline \textit{mmrt}.

Nếu đặt theo góc nhìn hội đồng chấm đồ án, một phiên bản được xem là đạt khi thỏa hai tầng bằng chứng tách biệt. Thứ nhất, baseline nghiên cứu phải có bằng chứng thực nghiệm trung thực trên dataset chuẩn. Thứ hai, application layer phải chạy ổn định ở mức demo và cho phép người xem theo dõi được luồng phân tích đầu-cuối. Cách định nghĩa này giúp báo cáo giữ đúng trọng tâm mô hình mà vẫn có phần triển khai đủ thuyết phục.

\section[Threats To Validity]{Threats to Validity}
Có bốn nhóm đe dọa chính tới độ vững của kết luận hiện tại. \textit{Construct validity} nằm ở biên ngữ nghĩa giữa \textit{looking\_away} và \textit{normal\_driving}: nếu độ lệch đầu hoặc ánh mắt chỉ xuất hiện ngắn và nhẹ, confusion pattern có thể phản ánh cả giới hạn mô hình lẫn giới hạn định nghĩa nhãn. \textit{Internal validity} nằm ở chỗ bundle \textit{v2} chưa đóng gói đủ log huấn luyện và ablation từng modality, nên vai trò của \textit{audio} hay \textit{skeleton} mới được diễn giải ở mức hợp lý kỹ thuật chứ chưa phải bằng chứng nhân quả.

\textit{External validity} bị giới hạn vì metric trung tâm hiện chỉ nằm trên validation của \textit{DBR-DMD-v1.1 CLEAN\_FULLMM\_5CLASS}, nên chưa thể suy rộng mạnh sang góc camera mới, người lái mới hay điều kiện ánh sáng ngoài tập đang có. Cuối cùng, \textit{deployment validity} nhắc rằng \textit{MDDAD} mới là tầng minh họa chuyển giao chứ chưa phải runtime thống nhất trùng hoàn toàn với baseline \textit{mmrt}; vì vậy kết quả vận hành của app không được đọc như benchmark cuối cùng của mô hình nghiên cứu.

\section[Han Che Va Huong Phat Trien]{Hạn Chế Và Hướng Phát Triển}
Báo cáo hiện còn bốn giới hạn trọng yếu: chưa xác minh \textit{test macro-F1} cho checkpoint tốt nhất; chưa có ablation từng modality cho chính \textit{v2}; chưa có benchmark chính quy về latency, throughput và realtime stability; và chưa đóng gói trọn bộ siêu tham số cuối cùng cùng log huấn luyện chi tiết. Nói gọn hơn, bản hiện tại mạnh hơn ở chiều \textit{analysis of available evidence} so với chiều \textit{complete experimental coverage}.

Từ đó, bốn hướng phát triển tự nhiên nhất là: (1) đóng gói lại toàn bộ train artifact của \textit{full\_v2}, gồm config, log và test summary; (2) thực hiện ablation theo modality và theo nhóm branch hình ảnh/phi hình ảnh; (3) thử các chiến lược hậu xử lý quyết định như smoothing, calibration và persistence rule cho cặp \textit{normal\_driving}/\textit{looking\_away}; và (4) tích hợp dần baseline nghiên cứu vào nhánh upload trước khi mở rộng sang realtime có benchmark tài nguyên rõ ràng.

\section[Ket Luan]{Kết Luận}
Báo cáo tập trung vào baseline nghiên cứu \textit{mmrt\_multimodal\_full\_v2\_t4x2} cho bài toán giám sát an toàn tài xế trong cabin. Mô hình này khai thác đồng thời \textit{body}, \textit{face}, \textit{audio} và \textit{skeleton}, đạt \textit{macro-F1 = 0.8066} trên tập validation của \textit{DBR-DMD-v1.1 CLEAN\_FULLMM\_5CLASS}.

Đóng góp chính của báo cáo không chỉ là con số \textit{0.8066}, mà là việc xác lập được một baseline đa phương thức thống nhất với biên bằng chứng rõ ràng: biết mô hình làm tốt ở đâu, yếu ở đâu, và phần nào đã được xác minh trực tiếp từ artifact. Insight trung tâm là độ khó không phân bố đều trên mọi lớp; \textit{phone\_use}, \textit{eye\_closure} và \textit{yawning} ổn định hơn, trong khi \textit{normal\_driving}/\textit{looking\_away} vẫn là biên quyết định khó nhất. Vì vậy, giá trị chính của \textit{mmrt\_multimodal\_full\_v2\_t4x2} nằm ở việc tạo ra một lõi nghiên cứu đủ mạnh để phân tích lỗi và đủ trung thực để không overclaim; còn \textit{MDDAD} chỉ giữ vai trò minh họa chuyển giao.

\begin{thebibliography}{10}
\small

\bibitem{mobilenetv3}
A. Howard, M. Sandler, G. Chu, et al.,
``Searching for MobileNetV3,''
\textit{Proceedings of ICCV}, 2019.

\bibitem{gru}
K. Cho, B. van Merrienboer, D. Bahdanau, and Y. Bengio,
``On the Properties of Neural Machine Translation: Encoder-Decoder Approaches,''
\textit{arXiv preprint arXiv:1409.1259}, 2014.

\bibitem{mediapipe}
C. Lugaresi, J. Tang, H. Nash, et al.,
``MediaPipe: A Framework for Building Perception Pipelines,''
\textit{arXiv preprint arXiv:1906.08172}, 2019.

\bibitem{dms}
S. Singh and N. Papanikolopoulos,
``Monitoring Driver Fatigue Using Facial Analysis Techniques,''
\textit{IEEE Intelligent Transportation Systems}, 1999.

\bibitem{slowfast}
C. Feichtenhofer, H. Fan, J. Malik, and K. He,
``SlowFast Networks for Video Recognition,''
\textit{Proceedings of ICCV}, 2019.

\bibitem{timesformer}
G. Bertasius, H. Wang, and L. Torresani,
``Is Space-Time Attention All You Need for Video Understanding?'',
\textit{Proceedings of ICML}, 2021.

\bibitem{videoswin}
Z. Liu, J. Ning, Y. Cao, et al.,
``Video Swin Transformer,''
\textit{Proceedings of CVPR}, 2022.

\bibitem{driveact}
M. Martin, A. Roitberg, M. Haurilet, and R. Stiefelhagen,
``Drive\&Act: A Multi-Modal Dataset for Fine-Grained Driver Behavior Recognition in Autonomous Vehicles,''
\textit{Proceedings of ICCV Workshops}, 2019.

\end{thebibliography}

\end{document}
